section{headline: Feedbackgespräche}

subheadline{Feedbackgespräch 1 10.03.2026 - Fragenkatalog}
Frageblöcke - Wichtige Aspekte --> Parameter sind zu definieren.

Nicht zu oberflächlich werden -> Zusammenhänge herausarbeiten!

"IT-Infrastruktur, die zur Anwendung kommt -> mobil, stationär, Gerät" -> Einfluss des Equipments auf die MFA

Welche Merkmale definieren technische Aspekte? Welche Aspekte definieren Regulatorik? -> jeweilige Richtlinien
Definition von IT-Governance (Reifegrad)

5 -> basierend worauf
5 -> weiß ich nicht
6 -> andere

Likert-Skala - nur die Grenzen
7-teilige Auswertung
Option: nicht relevant/keine Angabe

Größte Ablehnung = 1, Zustimmung = 7

Risiko abfragen ..? Kontrollfragen
Komplexität

8 -> Keine Dienste am Ende
Freiwilligkeit klären (FB, LI, IG, TT, WA)
-> LinkedIn als berufliche&private Plattform

Häufigsten Apps recherchieren -> diese explizit einbinden --> Auswahl begründen

nicht relevant = 0
Keine von den angeführten

Mehr Faktoren -> Komplexität abfragen

Zielvariable definieren -> 
Zuordnung vorbereiten -> Usability

Wie oft verwende ich MFA? -> Nutzungsfrequenz

Ziel: Verwendung quantitativ abbilden -> Kriterien aus der Literatur (zu klären)

Unterfrage 3 (Transfer klären) Was wird transferiert? -> Fokus "welche Merkmale der Technologie/welche Nutzungsfaktoren werden aus dem beruflichen Umgang in die private Verwendung der Technologie übernommen."

Unterfrage 1 erweitern ->  im beruflichen Kontext (Bankapp -> Zwang)

Zielgruppe definieren

Spill-over positiv/negativ? 
Frage, die den Wechsel abbildet --> wie stark motiviert mich das? --> Gegenfragen (Widerwille)
